% #############################################################################
% This is Chapter 5
% !TEX root = ../main.tex
% #############################################################################
\fancychapter{Experiments and Evaluation}
\label{chap:evaluation}

This chapter evaluates the trained models on variety identification and
bidirectional rewriting. \Cref{sec:evaluation-methodology} defines the
evaluation setup, and \Cref{sec:smaller-model-protocol-study} reports the protocol-selection
study with the 270M-parameter model used to choose the supervised training
protocol.
\Cref{sec:identification-results,sec:rewriting-results}
analyze identification and rewriting with the 4B-parameter model, while
\Cref{sec:error-analysis} examines copying behavior, representative
outputs, and limitations of the available supervision.

\section{Evaluation Setup}
\label{sec:evaluation-methodology}

Evaluation uses the FRMT test split and the Golden Collection described in
\Cref{chap:data}. Golden Collection evaluation uses the manually curated
reference instead of the automatic DeepL reference.

Rewriting quality is measured with BLEU \citep{papineni2002bleu}, computed with
SacreBLEU \citep{post-2018-call}, and TER \citep{snover2006study}, with higher
BLEU and lower TER indicating closer agreement with the reference. BLEU values use
the 0--100 reporting scale, and
rewriting uses beam search with beam size 4. Decoder-side variety labels are
removed before scoring. Rewriting comparisons are reported separately for
BP-to-EP and EP-to-BP when directional results are available, while the Stage A
ablation uses aggregate rewriting metrics. For evaluation, corpus BLEU and TER
are used for aggregate model comparisons, while sentence-level BLEU and TER are
reserved for the diagnostic analyses in \Cref{sec:error-analysis}. Stage C
separately uses sentence-level BLEU and TER within the training reward described
in \Cref{sec:method-stage-c}. Corpus BLEU is computed over
each complete dataset-and-direction slice by aggregating candidate \(n\)-gram
counts and the corresponding reference-limited match counts across all sentences
before calculating the corpus-level score.

Many EP/BP examples require only localized changes, and some require no
modification. The evaluation therefore also includes a deterministic copy
baseline that returns the source sentence unchanged, providing a reference for
how favorable each evaluation set is to source preservation.

Variety identification is evaluated with accuracy and macro-F1. Predictions are
obtained by scoring \texttt{PT} and \texttt{BR} under the adapted
sequence-to-sequence model and selecting the higher-scoring label.

\section[Supervised Protocol Selection with the 270M-Parameter T5Gemma 2 Model]{Supervised Protocol Selection with the\\270M-Parameter T5Gemma 2 Model}
\label{sec:smaller-model-protocol-study}

The supervised protocol was selected with the fully fine-tuned 270M-parameter
T5Gemma 2 model under decoder control so that alternative training choices
could be compared at lower computational cost. Table~\ref{tab:protocol-ablation-270m} reports the
base protocol, task-balanced batches, added-token treatment, and restricted
classification loss after the same OpenSubtitles Stage A and GPT-Wikipedia
Stage B path. The variants differ only in the protocol component defined in
\Cref{sec:method-ablation-design}.

\begin{table}[htbp]
\centering
\thesistablefont
\setlength{\tabcolsep}{3pt}
\caption{Decoder-controlled 270M full-fine-tuning supervised protocol-selection ablations.}
\label{tab:protocol-ablation-270m}
\begin{tabularx}{\textwidth}{@{}Y*{10}{c}@{}}
\toprule
\multirow{3}{=}{\textbf{Protocol variant}} & \multicolumn{5}{c}{\textbf{FRMT}} & \multicolumn{5}{c}{\textbf{Golden Collection}} \\
\cmidrule(lr){2-6} \cmidrule(lr){7-11}
& \multirow{2}{*}{\shortstack{\textbf{Macro-F1}\\\textbf{(\%)}}} & \multicolumn{2}{c}{\shortstack{\textbf{BLEU}\\\textbf{(0--100)}}} & \multicolumn{2}{c}{\textbf{TER}} & \multirow{2}{*}{\shortstack{\textbf{Macro-F1}\\\textbf{(\%)}}} & \multicolumn{2}{c}{\shortstack{\textbf{BLEU}\\\textbf{(0--100)}}} & \multicolumn{2}{c}{\textbf{TER}} \\
\cmidrule(lr){3-4} \cmidrule(lr){5-6} \cmidrule(lr){8-9} \cmidrule(lr){10-11}
& & {\scriptsize\textbf{BP-to-EP}} & {\scriptsize\textbf{EP-to-BP}} & {\scriptsize\textbf{BP-to-EP}} & {\scriptsize\textbf{EP-to-BP}} & & {\scriptsize\textbf{BP-to-EP}} & {\scriptsize\textbf{EP-to-BP}} & {\scriptsize\textbf{BP-to-EP}} & {\scriptsize\textbf{EP-to-BP}} \\
\midrule
Base protocol & 70.76 & 38.56 & 38.69 & \textbf{0.4748} & \textbf{0.4828} & 67.63 & 86.45 & 86.70 & 0.0978 & 0.0981 \\
Task-balanced batches & 69.74 & 38.56 & 38.55 & 0.4751 & 0.4859 & 68.52 & 86.04 & 86.45 & 0.0987 & 0.1001 \\
Added-token path & \textbf{72.93} & 38.36 & 38.67 & 0.4768 & 0.4864 & 66.66 & 86.23 & 86.60 & 0.0986 & 0.1006 \\
Restricted classification loss & 70.74 & \textbf{38.70} & \textbf{38.72} & 0.4750 & 0.4832 & \textbf{69.10} & \textbf{87.01} & \textbf{87.32} & \textbf{0.0941} & \textbf{0.0956} \\
\bottomrule
\end{tabularx}
\end{table}

Task-balanced batches do not provide a broad improvement: they slightly lower
FRMT identification and rewriting, improve Golden Collection identification,
and slightly lower Golden Collection rewriting. The added-token path gives the
largest FRMT identification gain, but
slightly lowers FRMT rewriting, lowers Golden Collection identification, and
does not improve Golden Collection rewriting. Restricted classification loss
gives the strongest Golden Collection identification and rewriting values, and
only a small FRMT BLEU increase, with a specialized classification loss but no
FRMT identification gain.

Although individual variants improve particular metrics, none provides a
consistent gain across identification and rewriting. The base protocol is
therefore retained as the simpler configuration for the subsequent experiments
with the 4B-parameter model. The retained protocol uses the natural
rewriting-to-identification mixture, the existing \texttt{PT}/\texttt{BR}
base-vocabulary representations, and normal decoder-vocabulary scoring at the
supervised identification label position, while rewriting uses loss over the
complete target sequence.

\section{Variety Identification Results}
\label{sec:identification-results}

Identification performance varies more clearly with the Stage B supervision
condition than with control placement, particularly on FRMT.
Table~\ref{tab:r48-classification-breakdown} reports these differences and the
Stage C results.

\begin{table}[htbp]
\centering
\thesistablefont
\setlength{\tabcolsep}{4pt}
\caption{Identification results for the 4B-parameter T5Gemma 2 model with rank-48 LoRA.}
\label{tab:r48-classification-breakdown}
\begin{tabularx}{\textwidth}{@{}clYcccc@{}}
\toprule
\textbf{Stage} & \textbf{Control} & \textbf{Data condition} &
\multicolumn{2}{c}{\textbf{FRMT}} &
\multicolumn{2}{c}{\textbf{Golden Collection}} \\
\cmidrule(lr){4-5} \cmidrule(lr){6-7}
& & & \shortstack{\textbf{Accuracy}\\\textbf{(\%)}} & \shortstack{\textbf{Macro-F1}\\\textbf{(\%)}} & \shortstack{\textbf{Accuracy}\\\textbf{(\%)}} & \shortstack{\textbf{Macro-F1}\\\textbf{(\%)}} \\
\midrule
B & Encoder & GPT-Wikipedia & 82.39 & 82.39 & \textbf{73.99} & \textbf{73.92} \\
B & Decoder & GPT-Wikipedia & 81.69 & 81.68 & 71.61 & 71.53 \\
B & Encoder & GPT-Wikipedia + FRMT & 86.09 & 86.06 & 71.43 & 71.19 \\
B & Decoder & GPT-Wikipedia + FRMT & \textbf{86.15} & \textbf{86.10} & 72.53 & 72.38 \\
\midrule
C & Encoder & GPT-Wikipedia & 82.33 & 82.33 & 73.63 & 73.50 \\
C & Decoder & GPT-Wikipedia & 81.52 & 81.31 & 72.34 & 71.81 \\
C & Encoder & GPT-Wikipedia + FRMT & 85.82 & 85.78 & 71.43 & 71.17 \\
C & Decoder & GPT-Wikipedia + FRMT & 85.65 & 85.54 & 72.34 & 72.02 \\
\bottomrule
\end{tabularx}
\end{table}

Adding FRMT supervision improves FRMT identification under both control
formulations, with macro-F1 increasing from 82.39 to 86.06 under encoder control
and from 81.68 to 86.10 under decoder control. The effect differs on the Golden
Collection: encoder-controlled macro-F1 decreases from 73.92 to 71.19, while
decoder-controlled macro-F1 increases from 71.53 to 72.38. Control placement
also changes with the supervision condition. With GPT-Wikipedia alone, encoder
control gives the higher macro-F1 on both evaluation resources, whereas after
adding FRMT, decoder control is slightly higher on both.

Accuracy and macro-F1 remain very close across all configurations, with
identification evaluated on equal numbers of PT- and BR-labeled examples.

The FRMT results also allow a comparison with the identification results
reported by \citet{sousa2025enhancingportuguesevarietyidentification}. The Stage
B encoder-controlled configuration trained on GPT-Wikipedia reaches 82.39
macro-F1 on FRMT, compared with the 77.25 F1 of their strongest
BERTimbau-based identifier. Although both studies use FRMT for evaluation,
their training setups and objectives differ: the previous model is dedicated
to identification, whereas this work uses a shared encoder-decoder interface.

Identification changes only slightly from Stage B to the corresponding Stage C
configurations. Across the four matched comparisons and both evaluation
resources, the largest absolute macro-F1 change is 0.56 percentage points, from
86.10 to 85.54 for the decoder-controlled GPT-Wikipedia + FRMT configuration on
FRMT. The remaining changes are smaller and include both increases and
decreases, with no consistent identification improvement or systematic
degradation across the evaluated conditions.

\section{Variety Rewriting Results}
\label{sec:rewriting-results}

Rewriting is more sensitive than identification to changes in supervision
composition. Table~\ref{tab:r48-translation-breakdown} summarizes performance
across supervision settings, control formulations, and training
stages.

\begin{table}[htbp]
\centering
\thesistablefont
\setlength{\tabcolsep}{3pt}
\caption{Directional rewriting results for the 4B-parameter T5Gemma 2 model after Stage B and Stage C with rank-48 LoRA.}
\label{tab:r48-translation-breakdown}
\begin{tabularx}{\textwidth}{@{}clYcccc@{}}
\toprule
\textbf{Stage} & \textbf{Control} & \textbf{Data condition} &
\multicolumn{2}{c}{\textbf{FRMT}} &
\multicolumn{2}{c}{\textbf{Golden Collection}} \\
\cmidrule(lr){4-5} \cmidrule(lr){6-7}
& & & \textbf{BP-to-EP} & \textbf{EP-to-BP}
& \textbf{BP-to-EP} & \textbf{EP-to-BP} \\
& & & \shortstack{\textbf{BLEU (0--100)}\\\textbf{/ TER}} & \shortstack{\textbf{BLEU (0--100)}\\\textbf{/ TER}}
& \shortstack{\textbf{BLEU (0--100)}\\\textbf{/ TER}} & \shortstack{\textbf{BLEU (0--100)}\\\textbf{/ TER}} \\
\midrule
B & Encoder & GPT-Wikipedia & 37.87/0.4791 & 40.10/0.4714 & 86.08/0.0946 & \textbf{89.09}/\textbf{0.0823} \\
B & Encoder (rewriting-only) & GPT-Wikipedia & 37.73/0.4803 & 40.31/0.4715 & 85.66/0.0976 & 88.66/0.0847 \\
B & Decoder & GPT-Wikipedia & 40.45/0.4594 & 40.60/0.4673 & 87.75/0.0874 & 87.50/0.0916 \\
B & Encoder & GPT-Wikipedia + FRMT & 37.99/0.4770 & \textbf{45.43}/0.4273 & 65.14/0.2555 & 66.61/0.2535 \\
B & Encoder (rewriting-only) & GPT-Wikipedia + FRMT & 36.76/0.4852 & 45.40/\textbf{0.4251} & 56.10/0.3222 & 58.69/0.3103 \\
B & Decoder & GPT-Wikipedia + FRMT & \textbf{43.97}/\textbf{0.4338} & 44.74/0.4306 & 68.37/0.2336 & 67.83/0.2400 \\
\midrule
C & Encoder & GPT-Wikipedia & 37.89/0.4784 & 40.28/0.4696 & 86.22/0.0953 & 88.93/0.0837 \\
C & Decoder & GPT-Wikipedia & 40.34/0.4603 & 40.33/0.4698 & \textbf{88.44}/\textbf{0.0824} & 88.02/0.0878 \\
C & Encoder & GPT-Wikipedia + FRMT & 38.35/0.4756 & 43.47/0.4440 & 77.16/0.1601 & 79.95/0.1488 \\
C & Decoder & GPT-Wikipedia + FRMT & 42.72/0.4417 & 42.72/0.4466 & 81.52/0.1351 & 80.70/0.1402 \\
\bottomrule
\end{tabularx}
\end{table}

Adding FRMT supervision has a direction-dependent effect on FRMT rewriting,
ranging from almost no change to a clear improvement. In contrast, Golden
Collection BLEU drops by roughly 20 points across every direction and control
comparison, with corresponding increases in TER.

To examine whether identification supervision changes rewriting quality, the
encoder-controlled Stage B comparison includes matched multitask and
rewriting-only configurations trained with the same rewriting data. With
GPT-Wikipedia, removing identification supervision changes rewriting by less
than half a BLEU point in every comparison. With GPT-Wikipedia + FRMT, the two
configurations remain close on FRMT, while the multitask configuration is
stronger on the Golden Collection by 9.04 BLEU in BP-to-EP and 7.92 BLEU in
EP-to-BP. Removing identification supervision does not improve
encoder-controlled rewriting.

The effect of control placement depends on rewriting direction. Decoder control
performs better in every BP-to-EP comparison, while EP-to-BP is mixed and
sometimes favors encoder control.

Beyond these supervised comparisons, Stage C shows how the models change after
reward-based continuation. Reward-based continuation changes the GPT-Wikipedia models only slightly. The
larger effect appears after GPT-Wikipedia + FRMT specialization, where Golden
Collection BLEU recovers by roughly 12--13 points across the direction and
control comparisons. Changes in FRMT rewriting performance remain comparatively
small and mixed. Stage C therefore recovers much of the Golden Collection
performance lost after Stage B specialization with the GPT-Wikipedia + FRMT data
mix, while its FRMT effect remains direction- and formulation-dependent.

The most directly comparable prior rewriting result is BP-to-EP evaluation on
the Golden Collection. After rescaling the BLEU values reported by
\citet{sanches2024brazilianportugueseeuropeanportuguese} to the 0--100 scale,
their strongest trained system obtains approximately 84 BLEU and 0.10 TER,
while their published single-direction BP-to-EP copy baseline reaches
approximately 89 BLEU and 0.07 TER. The decoder-controlled GPT-Wikipedia
configuration reaches 87.75 BLEU and 0.0874 TER at Stage B and 88.44 BLEU and
0.0824 TER after Stage C. This numerical relation is contextual, as
\citet{sanches2024brazilianportugueseeuropeanportuguese} train specifically for
BP-to-EP rewriting, while the present configuration is bidirectional and also
supports identification.

Additional descriptive comparisons with external pretrained and
instruction-following models evaluated under the same protocol are reported in
\Cref{chap:appendix-external-model-comparison}.

The Stage A ablation reported in Table~\ref{tab:stage-a-ablation} compares each
configuration with a matched model that starts directly from Stage B. To capture
effects on both tasks, the comparison reports macro-F1 together with aggregate
rewriting BLEU and TER deltas.

\begin{table}[htbp]
\centering
\thesistablefont
\setlength{\tabcolsep}{4pt}
\caption{Stage A warm-up ablation under aggregate evaluation. Deltas are computed as the Stage A model minus the matched no-Stage-A model. Positive deltas indicate improvement for macro-F1 and BLEU, while negative deltas indicate improvement for TER.}
\label{tab:stage-a-ablation}
\begin{tabularx}{\textwidth}{@{}llYccc@{}}
\toprule
\textbf{Control} & \textbf{Data condition} & \textbf{Dataset} & \shortstack{\(\Delta\) \textbf{Macro-F1}\\\textbf{(pp)}} & \shortstack{\(\Delta\) \textbf{BLEU}\\\textbf{(points)}} & \(\Delta\) \textbf{TER} \\
\midrule
Encoder & GPT-Wikipedia & FRMT & \textbf{+0.87} & \textbf{+4.40} & \textbf{-0.0573} \\
Encoder & GPT-Wikipedia & Golden Collection & \textbf{+4.01} & \textbf{+25.41} & \textbf{-0.2404} \\
Encoder & GPT-Wikipedia + FRMT & FRMT & +0.81 & +1.92 & -0.0176 \\
Encoder & GPT-Wikipedia + FRMT & Golden Collection & +1.62 & -6.79 & +0.0425 \\
Decoder & GPT-Wikipedia & FRMT & -1.28 & -0.01 & +0.0004 \\
Decoder & GPT-Wikipedia & Golden Collection & +2.67 & +1.42 & -0.0144 \\
Decoder & GPT-Wikipedia + FRMT & FRMT & -0.07 & +0.35 & -0.0031 \\
Decoder & GPT-Wikipedia + FRMT & Golden Collection & +2.19 & +2.00 & -0.0217 \\
\bottomrule
\end{tabularx}
\end{table}

The clearest aggregate Stage A benefit occurs on the Golden Collection
after specialization with GPT-Wikipedia data, especially under encoder control,
where Stage A adds 4.01 macro-F1 points and 25.41 BLEU while reducing TER by
0.2404. The decoder-controlled GPT-Wikipedia condition also improves on the
Golden Collection, while the GPT-Wikipedia + FRMT conditions show smaller or
mixed changes.

Overall, supervision composition produces the clearest changes in rewriting
performance. GPT-Wikipedia alone is consistently better suited to the Golden
Collection at Stage B, while adding FRMT can improve FRMT rewriting depending
on direction and control formulation but sharply reduces Golden Collection
performance. Decoder control shows the stronger overall tendency for rewriting,
performing better in every BP-to-EP comparison, although EP-to-BP remains mixed.
Stage A generally helps the training process, with its largest gains after
GPT-Wikipedia specialization, while Stage C produces its largest gains after
GPT-Wikipedia + FRMT specialization, where it consistently recovers Golden
Collection performance across control formulations and rewriting directions.

\section{Error Analysis}
\label{sec:error-analysis}

\Cref{sec:rewriting-results} shows that adding FRMT supervision can improve FRMT
rewriting while sharply reducing Golden Collection performance at Stage B. To
examine this contrast, \Cref{sec:error-analysis-copy} first considers how closely
references remain to their sources,
\Cref{sec:error-analysis-frmt-stage-b} then compares the copying and rewriting
behavior under the two Stage B conditions, and
\Cref{sec:error-analysis-data} examines supervision limitations.

\subsection{Copy-Friendly Evaluation Cases}
\label{sec:error-analysis-copy}

The Golden Collection contains many cases in which little or no rewriting is
required. Of its 1,000 directional cases, 454 have identical source and
reference sentences. When the source itself is used as the candidate, 566 cases
obtain TER at most 0.05 and 636 obtain TER at most 0.08. At corpus level, the
unchanged-source baseline reaches 89.17 BLEU and 0.0755 TER on the Golden
Collection, compared with 38.75 BLEU and 0.4819 TER on FRMT. The Golden
Collection is therefore much more favorable to conservative outputs, whereas
FRMT references differ more substantially from their sources.

\Cref{tab:golden-copy-comparison-rates} compares sentence-level BLEU (sBLEU)
and TER with the unchanged-source baseline.

\begin{table}[htbp]
\centering
\thesistablefont
\setlength{\tabcolsep}{4pt}
\caption{Golden Collection sentence-level predictions matching or outperforming the unchanged-source baseline for the Stage B models in Table~\ref{tab:r48-translation-breakdown}.}
\label{tab:golden-copy-comparison-rates}
\begin{tabularx}{\textwidth}{@{}Ylccc@{}}
\toprule
\textbf{Data condition} &
\textbf{Control} &
\textbf{Direction} &
\shortstack{\textbf{sBLEU \(\geq\) copy}\\\textbf{(\%)}} &
\shortstack{\textbf{TER \(\leq\) copy}\\\textbf{(\%)}} \\
\midrule
GPT-Wikipedia & Encoder & BP-to-EP & 78.6 & 78.8 \\
GPT-Wikipedia & Encoder & EP-to-BP & 81.4 & 81.2 \\
GPT-Wikipedia & Decoder & BP-to-EP & 74.0 & 74.4 \\
GPT-Wikipedia & Decoder & EP-to-BP & 72.8 & 73.6 \\
\midrule
GPT-Wikipedia + FRMT & Encoder & BP-to-EP & 23.0 & 23.0 \\
GPT-Wikipedia + FRMT & Encoder & EP-to-BP & 23.6 & 24.4 \\
GPT-Wikipedia + FRMT & Decoder & BP-to-EP & 31.0 & 31.2 \\
GPT-Wikipedia + FRMT & Decoder & EP-to-BP & 30.2 & 29.8 \\
\bottomrule
\end{tabularx}
\end{table}

The sentence-level comparison in
\Cref{tab:golden-copy-comparison-rates} counts a prediction when its
sentence-level BLEU matches or exceeds the copy baseline, or when its
sentence-level TER matches or improves on the copy baseline. Across both control
formulations, predictions under GPT-Wikipedia supervision match or outperform
copying in 72.8--81.4\% of cases by sentence-level BLEU and 73.6--81.2\% by
sentence-level TER. With GPT-Wikipedia + FRMT, the corresponding ranges fall to
23.0--31.0\% and 23.0--31.2\%, respectively.

Sentence-level comparisons cannot be read as sentence-weighted versions of the
corpus metrics. Corpus BLEU aggregates \(n\)-gram match statistics over the complete
evaluation slice instead of averaging sentence-level BLEU or counting wins. A
model can therefore match or improve over copying on a large proportion of
individual sentences yet obtain lower corpus BLEU when losses on other examples remove
proportionally more matched \(n\)-grams. This sensitivity is especially relevant
for the many short, copy-friendly Golden Collection sentences. Changing one word
in a short sentence can disrupt a large fraction of its available \(n\)-grams,
so an unnecessary change can incur a disproportionate overlap loss. A correct
small edit and an unnecessary small edit can consequently have very different
effects depending on the reference. Corpus TER similarly sums edit counts over
the complete evaluation slice and divides by the total reference-token count. It
does not average sentence-level TER values, so sentence-level gains
can coexist with worse corpus TER when losses elsewhere add more edits
per reference token.

\Cref{tab:golden-copy-bleu-paradox} illustrates both outcomes with one
conservative Stage B GPT-Wikipedia checkpoint. In the first and second rows, the
model applies the required local substitution and matches the reference exactly.
In the third and fourth rows, copying is correct, but rewriting sharply lowers
sentence-level BLEU.

\begin{table}[htbp]
\centering
\thesistablefont
\setlength{\tabcolsep}{4pt}
\caption{Golden Collection examples comparing the unchanged-source baseline with predictions from a conservative Stage B GPT-Wikipedia checkpoint. Boldface marks the sentence spans responsible for the illustrated contrast, and the higher sBLEU value in each row is bold.}
\label{tab:golden-copy-bleu-paradox}
\begin{tabularx}{\textwidth}{@{}YYYcc@{}}
\toprule
Source & Reference & Model prediction & \shortstack{Copy sBLEU\\(0--100)} & \shortstack{Model sBLEU\\(0--100)} \\
\midrule
\textit{Clonar animais continua \textbf{sendo} uma tarefa dif\'icil.} &
\textit{Clonar animais continua \textbf{a ser} uma tarefa dif\'icil.} &
\textit{Clonar animais continua \textbf{a ser} uma tarefa dif\'icil.} &
44.12 &
\textbf{100.00} \\

\textit{O artigo do \textbf{investigador} est\'a na edi\c{c}\~ao de hoje da revista ``Science''.} &
\textit{O artigo do \textbf{pesquisador} est\'a na edi\c{c}\~ao de hoje da revista ``Science''.} &
\textit{O artigo do \textbf{pesquisador} est\'a na edi\c{c}\~ao de hoje da revista ``Science''.} &
80.03 &
\textbf{100.00} \\

\textit{\textbf{V.Ex.\textordfeminine}\ est\'a inscrito para falar contra.} &
\textit{\textbf{V.Ex.\textordfeminine}\ est\'a inscrito para falar contra.} &
\textit{\textbf{Sua Excel\^encia} est\'a inscrito para falar contra.} &
\textbf{100.00} &
46.76 \\

\textit{--- Ah!} &
\textit{--- Ah!} &
\textit{\textbf{\'E...}} &
\textbf{100.00} &
0.00 \\
\bottomrule
\end{tabularx}
\end{table}

The much weaker FRMT copy baseline reflects references that are farther from
their sources on average. Broader rewriting can therefore improve corpus-level
agreement more readily than on the Golden Collection, although broader changes
are not necessarily appropriate for every example in FRMT.

\subsection{Effect of FRMT During Supervised Specialization}
\label{sec:error-analysis-frmt-stage-b}

Given the Golden Collection's sensitivity to unnecessary changes, the next
question is whether the two Stage B supervision conditions differ in how often
they preserve the source. The encoder-controlled Stage B models trained with
GPT-Wikipedia data and with the GPT-Wikipedia + FRMT data mix differ clearly in
this behavior. On the 454 exact-copy Golden Collection rows, where leaving the source
unchanged matches the reference, the model trained with GPT-Wikipedia preserves
the source exactly in 74.7\% of cases, compared with 18.7\% for the model trained
with the GPT-Wikipedia + FRMT data mix. On the 546 non-exact rows, the
corresponding copy rates are 33.7\% and 3.1\%. In these cases, unchanged copying
can indicate under-editing, so reducing copying can help realize required
changes but also increases the opportunity for unnecessary reformulation.

\Cref{tab:golden-qualitative-overtranslation} illustrates both sides of this
trade-off. In the first row, the model trained with GPT-Wikipedia under-edits by
leaving \textit{procura} unchanged, whereas the model trained with the
GPT-Wikipedia + FRMT data mix makes the required change to \textit{demanda} but
also reformulates the sentence more broadly. In the fourth row, the reference is
identical to the source, which the GPT-Wikipedia model preserves, while the
GPT-Wikipedia + FRMT model rewrites unnecessarily.

\begin{table}[htbp]
\centering
\thesistablefont
\setlength{\tabcolsep}{4pt}
\caption{Illustrative Golden Collection predictions from encoder-controlled Stage B models trained with GPT-Wikipedia and GPT-Wikipedia + FRMT. Boldface marks model-output spans that illustrate under-editing or unnecessary reformulation discussed in the text.}
\label{tab:golden-qualitative-overtranslation}
\begin{tabularx}{\textwidth}{@{}cYYYY@{}}
\toprule
Example & Source & Reference & GPT-Wikipedia & GPT-Wikipedia + FRMT \\
\midrule
1 &
\textit{``A procura aumenta de forma dr\'astica no mundo todo'', afirmou o especialista em recursos h\'idricos Jos\'e Galizia Tundisi, do Instituto Internacional de Ecologia, em S\~ao Carlos (SP).} &
\textit{``A demanda aumenta de forma dr\'astica no mundo todo'', afirmou o especialista em recursos h\'idricos Jos\'e Galizia Tundisi, do Instituto Internacional de Ecologia, em S\~ao Carlos (SP).} &
\textit{``A \textbf{procura} aumenta de forma dr\'astica no mundo todo'', afirmou o especialista em recursos h\'idricos Jos\'e Galizia Tundisi, do Instituto Internacional de Ecologia, em S\~ao Carlos (SP).} &
\textit{``A demanda \textbf{est\'a aumentando drasticamente em todo o mundo}'', \textbf{disse} o especialista em recursos h\'idricos Jos\'e Galizia Tundisi, do Instituto Internacional de Ecologia em S\~ao Carlos (SP).} \\

2 &
\textit{Ele afirmou que o aumento de temperatura, que pode chegar a 3C nas \'areas h\'umidas, seria um desastre.} &
\textit{Ele afirmou que o aumento de temperatura, que pode chegar a 3C nas \'areas \'umidas, seria um desastre.} &
\textit{Ele afirmou que o aumento de temperatura, que pode chegar a 3C nas \'areas \'umidas, seria um desastre.} &
\textit{Ele \textbf{disse} que o aumento de temperatura, que pode chegar a \textbf{3\textordmasculine C em \'areas \'umidas}, seria um desastre.} \\

3 &
\textit{Segundo o investigador, a equipa inclusive pediu patentes sobre algumas dessas aplica\c{c}\~oes, mas, enquanto o pedido n\~ao for aprovado, Perales prefere n\~ao revelar exatamente do que se trata.} &
\textit{Segundo o pesquisador, a equipe inclusive pediu patentes sobre algumas dessas aplica\c{c}\~oes, mas, enquanto o pedido n\~ao for aprovado, Perales prefere n\~ao revelar exatamente do que se trata.} &
\textit{Segundo o pesquisador, a equipe inclusive pediu patentes sobre algumas dessas aplica\c{c}\~oes, mas, enquanto o pedido n\~ao for aprovado, Perales prefere n\~ao revelar exatamente do que se trata.} &
\textit{Segundo o pesquisador, a equipe \textbf{at\'e mesmo solicitou} patentes sobre algumas dessas aplica\c{c}\~oes, mas enquanto o pedido n\~ao for aprovado, Perales prefere n\~ao revelar exatamente do que se trata.} \\

4 &
\textit{Da\'i por que est\'a absolutamente correto no conte\'udo, como na forma, o parecer do ilustre Relator Bernardo Cabral.} &
\textit{Da\'i por que est\'a absolutamente correto no conte\'udo, como na forma, o parecer do ilustre Relator Bernardo Cabral.} &
\textit{Da\'i por que est\'a absolutamente correto no conte\'udo, como na forma, o parecer do ilustre Relator Bernardo Cabral.} &
\textit{\textbf{Por isso, o parecer do ilustre Relator, Bernardo Cabral, \'e absolutamente correto tanto no conte\'udo como na forma.}} \\
\bottomrule
\end{tabularx}
\end{table}

\subsection{Limitations of the Available Supervision}
\label{sec:error-analysis-data}

The broader rewriting behavior observed after adding FRMT supervision may
reflect differences in the aligned data itself. Inspection of the available
pairs shows source--reference differences that extend beyond localized EP/BP
conversion, including lexical reformulation and changes in expressed
information. In some FRMT examples, lexical substitutions such as
\textit{marcante} to \textit{importante},
\textit{ganhou} to \textit{recebeu},
\textit{reside} to \textit{vive},
\textit{usar} to \textit{utilizar}, and
\textit{m\'usica} to \textit{can\c{c}\~ao} alter the wording
without a clear EP/BP adaptation.

A second issue concerns differences in the information expressed by the source
and reference. \Cref{tab:dataset-information-mismatch} shows representative
cases from both evaluation resources. In the Golden Collection examples, the
reference removes or compresses source material, while in the FRMT examples it
removes repeated title glosses and unit conversions. These transformations go
beyond direct language-variety substitution. Additional inspected FRMT examples
contain other omissions or editorial changes of the same general kind, so the
cases in \Cref{tab:dataset-information-mismatch} are illustrative rather than
exhaustive.

\begin{table}[H]
\centering
\thesistablefont
\setlength{\tabcolsep}{4pt}
\caption{Examples where source and reference sentences differ in expressed information. Boldface marks the spans responsible for the information difference discussed in the text.}
\label{tab:dataset-information-mismatch}
\begin{tabularx}{\textwidth}{@{}lcYY@{}}
\toprule
Dataset & Direction & Source & Reference \\
\midrule
Golden Collection &
BP-to-EP &
\textit{- Oww voc\^e \'e um amor Obrigado viu!- Falei \textbf{e o mesmo retribuiu com um abra\c{c}o.}} &
\textit{- Oww, \'es ador\'avel. Obrigado!} \\

Golden Collection &
EP-to-BP &
\textit{Ent\~ao, o latifundi\'ario pode tirar toda a madeira, expulsar os \textbf{que ocuparam a terra abandonada} e fazer verdadeiros \textbf{atos} absurdos nesse per\'iodo.} &
\textit{Ent\~ao, o latifundi\'ario pode tirar toda a madeira, expulsar os \textbf{posseiros} e fazer verdadeiros absurdos nesse per\'iodo.} \\

FRMT &
BP-to-EP &
\textit{A UNICEF concedeu a Lerner o Pr\^emio Crian\c{c}a e Paz em 1996 pelos seus programas ``Da Rua para a Escola'' \textbf{(Da Rua para a Escola)}, ``Protegendo a Vida'' \textbf{(Protegendo a Vida)} e ``Universidade do Professor'' \textbf{(Universidade do Professor)}.} &
\textit{A UNICEF atribuiu a Lerner o Pr\'emio Crian\c{c}a e Paz em 1996 pelos programas ``Da Rua para a Escola'', ``Protegendo a Vida'' e ``Universidade do Professor''.} \\

FRMT &
BP-to-EP &
\textit{O Aut\'odromo Internacional Orlando Moura tem uma pista de 3.433 metros \textbf{(11.263 p\'es)}, e o Kart\'odromo Ayrton Senna uma pista de 930 metros \textbf{(3.051 p\'es)}.} &
\textit{O Aut\'odromo Internacional Orlando Moura tem uma pista de 3433 metros e o Kart\'odromo Ayrton Senna tem uma pista de 930 metros.} \\
\bottomrule
\end{tabularx}
\end{table}

Such differences affect both training and evaluation. When they appear in
supervision, they can teach transformations broader than controlled EP/BP
rewriting. When they appear in evaluation references, they can affect scores
even for meaning-preserving, appropriate variety-specific changes.

This analysis also highlights a limitation of BLEU and TER. Both compare a
model output with a particular reference, so they cannot determine whether a
difference from that reference reflects a necessary variety-specific edit, an
acceptable alternative, or an unnecessary reformulation. This ambiguity
becomes more important when the reference itself contains paraphrastic or
editorial changes. The same issue also carries into Stage C, where
sentence-level BLEU and TER contribute directly to the reward. If a reference
contains a broader rewrite than the task requires, the reward can favor
candidates that reproduce that rewrite over candidates that make a smaller but
still appropriate variety-specific change. Conversely, when copying is
appropriate, the relative-to-copy formulation can reward restraint. The
reference quality therefore affects both evaluation and the Stage C training
signal.
